
\begingroup
\shorthandoff{<>}% Babel español trata de otro modo ciertos delimitadores de Xy-pic.
\begin{corollaire}
\label{XIII.2.5}
Sean
\marginpar{385}
$f\colon X\to S$ un morfismo, $D$ un divisor sobre~$X$
de cruzamientos normales relativamente a~$S$ \eqref{XIII.2.1},
$Y=\Supp D$, $U=X-Y$, e $i\colon U\to X$ la inmersión
canónica. Sea $\Phi$ un stack sobre~$U$ y supongamos dados
morfismos étale sobreyectivos $S_1\to S$ y $X_2\to
X\times_{S}S_1$, un stack $\Psi$ sobre~$S_1$ y un isomorfismo
$\Phi|U_2\simeq\Psi|U_2$ (\cf \Ref{XIII.2.2}).

Entonces, para todo morfismo $h\colon S'\to S$, si $k=h_{(X)}$,
el funtor canónico
$$
\varphi\colon k^*i_*\Phi\to i_*'\Phi'
$$
es plenamente fiel. Si $\Psi$ es constructible, el funtor
canónico
$$
\psi\colon k^*i_*^{\tame}\Phi\to i_*^{\prime\tame}\Phi'
$$
es una equivalencia de categorías.

Además, si el stack $\Psi$ es constructible (\resp si $\Psi$ es
$1$-constructible \eqref{XIII.0} y $S$ localmente noetheriano),
$i_*\Phi$ es constructible (\resp
$i_*^t\Phi$
es $1$-constructible).
\end{corollaire}

Demostremos que $\varphi$ es plenamente fiel. Basta ver
que, para todo punto geométrico $\overline{x}'$ de~$X'$, lo
mismo ocurre con el funtor
$$
\overline{\varphi}\colon
\Phi(\overline{U})\to\Phi'(\overline{U'})
$$
(se retoman las notaciones de~\Ref{XIII.2.4.2}). Sean $a$, $b$ dos
objetos de~$\Phi(\overline{U})$, y sean $a',\ b'$ sus imágenes por
$\overline{\varphi}$. Como el morfismo
$\overline{U}\to\overline{S}$ es localmente $0$-acíclico
\eqref{XIII.5.7}, se tiene un isomorfismo
$$
\H^0(\overline{U},\overline{S\Phi})
\overset{\sim}{\to}
\ifthenelse{\boolean{orig}}
{\H^0(S,\overline{S\Psi})}
{\H^0(\overline{S},\overline{S\Psi})}.
$$
Por consiguiente, $a$ y $b$ provienen por imagen inversa de objetos
de~$\Psi$, y lo mismo ocurre, por tanto, con
$F=\SheafHom_{\overline{U}}(a,b)$. Como el haz
$F'=\SheafHom_{\overline{U'}}(a',b')$ es la imagen inversa
sobre~$\overline{U'}$ de~$F$, de~\Ref{XIII.2.4.1} resulta que el
morfismo canónico
$$
\H^0(\overline{U},F)\to \H^0(\overline{U'},F')
$$
es un isomorfismo, lo que demuestra que $\overline{\varphi}$ es
plenamente fiel.

Demostremos
\marginpar{386}
que, para todo punto geométrico $\overline{x}'$ de~$X'$, el
funtor
$$
\overline{\psi}\colon
i^{\tame}_*\Phi(\overline{X'})
\to i_*^{\prime\tame}\Phi'(\overline{X'})
$$
es una equivalencia. Por lo anterior,
$\overline{\psi}$ es plenamente fiel. Demostremos que
$\overline{\psi}$ es esencialmente sobreyectivo. Sea $a'$ un
objeto de~$\Phi'(\overline{U'})$ moderadamente
ramificado sobre~$X'$ relativamente a~$S'$ \eqref{XIII.2.2}, y
demostremos que es la imagen de un objeto moderadamente
ramificado de~$\Phi(\overline{U})$. De~\Ref{XIII.2.4}
1) resulta que el morfismo canónico
$$
\H^0(\overline{U},\overline{S}\overline{\Phi})\to
\H^0(\overline{U'},\overline{S}\overline{\Phi}')
$$
es un isomorfismo. Sea $G'$ el subgerbe maximal de~$\Phi'$
generado por~$a'$; existe entonces un subgerbe maximal~$G$ de
$\overline{\Phi}$, imagen inversa de un gerbe sobre~$\overline{S}$, tal
que se tiene
$$
\overline{m}^*G\simeq G',
$$
donde
\ifthenelse{\boolean{orig}}
{$m$}
{$\overline{m}$}
es el morfismo $\overline{U'}\to
\overline{U}$. El funtor canónico
\begin{equation*}
\label{eq:XIII.2.5.*}
\tag{$*$} {\overline{k}^*\overline{i}_*^{\tame}G\to
\overline{i}_*^{\prime\tame}G'}
\end{equation*}
es una equivalencia, pues $G$ se identifica con un gerbe de
torsores bajo un haz de grupos constructible proveniente de
$\overline{S}$, y puede aplicarse \Ref{XIII.2.4}~1). De
\eqref{eq:XIII.2.5.*} resulta entonces que existe un
objeto $a$ de~$G(\overline{U})$, moderadamente
ramificado sobre~$X$ relativamente a~$S$, cuya imagen inversa sobre
$\overline{U'}$ es $a'$, lo que demuestra que $\psi$ es una
equivalencia.

Si $\Psi$ es constructible, lo mismo ocurre con~$i_*\Phi$; en
efecto, un objeto $x$ de~$i_*\Phi$ es, localmente para la topología
étale sobre~$S$, imagen inversa de un objeto $y$ de~$\Psi$; de
\Ref{eq:XIII.2.4.1.1} resulta, por tanto, que $\SheafAut (x)$ es
la imagen inversa de~$\SheafAut (y)$, luego es constructible. Finalmente, si
$\Psi$ es $1$-constructible, lo mismo ocurre con
$i^{\tame}_*\Phi$ por~\Ref{XIII.6.3} más abajo.

\begin{corollaire}
\label{XIII.2.6}
Las
\marginpar{387}
notaciones son las de~\Ref{XIII.2.4}. Supongamos que $S$ tiene
característica cero en todo punto $s$ tal que $Y_s\neq
\emptyset.$ Entonces, si $\cal{ F}$ es un haz de grupos localmente
constante y constructible sobre~$U$ (\resp un stack constructible
sobre~$U$ que, localmente sobre~$X$ y~$S$, es imagen inversa de un
stack constructible sobre~$S$), $(\cal{ F},i)$ es cohomológicamente
propio relativamente a~$S$ en dimensión $\leq 1.$
\end{corollaire}

Como todo haz de conjuntos constructible sobre~$U$ está
moderadamente ramificado sobre~$X$ relativamente a~$S,$ el
corolario resulta de~\Ref{XIII.2.4} (\resp \Ref{XIII.2.5}).

\begin{corollaire}
\label{XIII.2.7}
Las notaciones son las de~\Ref{XIII.2.4}, pero se da además
un $S$-esquema $T$, un morfismo propio $p\colon X\to T,$
y se supone que $X$ y $T$ son de presentación finita sobre~$S$; sea
$q=pi$. Sea $\cal{ F}$ un haz de conjuntos constructible sobre~$U$
que satisfaga una de las condiciones \textup{a)} o \textup{b)}
de~\Ref{XIII.2.4} (\resp un haz de grupos que satisfaga una de las
condiciones \textup{a)} o \textup{b)} de~\Ref{XIII.2.4},
\resp un stack sobre~$U$ que, localmente sobre~$X$ y~$S$, es imagen
inversa de un stack constructible $G$ sobre~$S$). Entonces
se tienen las conclusiones siguientes:
\begin{enumerate}
\item[1)] $(\cal{ F},q)$ es cohomológicamente propio relativamente
a~$S$ en dimensión $\leq 0$ (\resp para todo morfismo $h\colon
S^{'}\to S$, si $m=h_{(T)}$, el morfismo canónico
$$
\Theta\colon m^*(\R^1_{\tame}q_*\cal{ F})\to
\R^1_{\tame}q_*^{'}\cal{ F}^{'}
$$
es un isomorfismo, \resp
\ifthenelse{\boolean{orig}}
{para todo morfismo $S^{'}\to S$,}
{para todo morfismo $h\colon S^{'}\to S$, si $m=h_{(T)}$,}
el morfismo canónico
$$
\xi \colon m^*(q_*^{\tame}\cal{ F})\to
{q'}_*^{{\tame}}\cal{F}^{'}
$$
es una equivalencia).
\item[2)] El haz $q_*\cal{ F}$ (\resp el haz
$\R_{\tame}^1q_*\cal{ F}$, \resp el stack $q_*^{\tame}\cal{ F}$)
es constructible. En el último caso, si se supone que $S$ es localmente
noetheriano y que $G$ es $1$-constructible, lo mismo ocurre con
$q_*^{\tame}\cal{ F}.$
\end{enumerate}
\end{corollaire}

La
\marginpar{388}
primera parte resulta inmediatamente de~\Ref{XIII.2.4},
\Ref{XIII.2.5} y de la demostración
de~\Ref{XIII.1.8}. Demostremos~2). Si $\cal{ F}$ es un haz
de conjuntos constructible sobre~$U$ que satisface \Ref{XIII.2.4}~a)
o \Ref{XIII.2.4}~b), de \Ref{XIII.2.4}~2) resulta que
$i_*\cal{ F}$ es constructible; lo mismo ocurre, por tanto, con
$q_*\cal{ F}=p_*(i_*\cal{ F})$ (SGA~4 XIV~1.1).

Sea $\cal{F}$ un haz de grupos constructible sobre~$U$
que satisfaga \Ref{XIII.2.4}~a) o \Ref{XIII.2.4}~b), y demostremos
que $\R_{\tame}^1q_*\cal{ F}$ es constructible. Por paso al
límite (EGA~IV 8.10.5 y 17.7.8) y usando~1), puede
suponerse que $S$ es noetheriano. Sea entonces $\Phi$ el stack sobre~$X$
cuya fibra sobre todo esquema $X^{'}$ étale sobre~$X$ está formada
por los torsores sobre~$U^{'}=U\times_{X}X^{'}$, de grupo
\ifthenelse{\boolean{orig}}
{$\cal{F}|U$}
{$\cal{F}|U'$},
que están moderadamente ramificados sobre
\ifthenelse{\boolean{orig}}
{~$X$}
{~$X'$}
relativamente a~$S$; se tiene, por tanto,
$$
\ifthenelse{\boolean{orig}}
{S(i_*^{\tame}\Phi)}
{S(\Phi)}
\simeq \R_{\tame}^1i_*\cal{ F}
$$
y este haz es constructible por~\Ref{XIII.2.4}~2).
De ello resulta, por
\ifthenelse{\boolean{orig}}
{\Ref{XIII.6.3}}
{\Ref{XIII.6.2}}
más abajo, que $S(p_*\Phi)$
es constructible, \ie que $\R^1_{\tame}q_*\cal{ F}$ es
constructible.

Finalmente, si $\cal{ F}$ es un stack sobre~$U$ que, localmente
sobre~$X$ y $S$, es imagen inversa de un stack constructible sobre~$S,$
$i_*^{\tame}\cal{ F}$ es constructible, y lo mismo ocurre, por tanto,
con $q_*^{\tame}\cal{ F}=p_*(i_*^{\tame}\cal{ F})$. Si
además $S$ es localmente noetheriano y $SG$ es constructible,
$S(i_*^{\tame}\cal{ F})$ es constructible por~6.3; lo mismo ocurre,
por tanto, con
\ifthenelse{\boolean{orig}}
{$S(q_*i_*^{\tame}\cal{ F})$}
{$S(p_*i_*^{\tame}\cal{ F})$},
\ie con
$S(q_*^{\tame}\cal{ F})$~\Ref{XIII.6.2}.

\begin{corollaire}
\label{XIII.2.8}
Sea
$$
\xymatrix{ U \ar[r]^{i} \ar[d]_{g} & X \ar[ld]^{f} \\ S }
$$
un diagrama conmutativo de esquemas, en el que $U$ es el abierto
complementario
\marginpar{389}
en $X$ de un divisor de cruzamientos normales relativamente a~$S$, y
$f$ un morfismo \emph{propio} de presentación finita. Sea
$\LL$ un conjunto de números primos. Supongamos que~$g$ es localmente
$0$-acíclico (\resp localmente $1$-asférico para $\LL$). Entonces, si
$\cal{ F}$ es un haz de conjuntos sobre~$U$ (\resp un haz de $\LL$-grupos)
sobre~$U,$ localmente constante y constructible,
moderadamente ramificado sobre~$X$ relativamente a~$S$,
\ifthenelse{\boolean{orig}}
{$f_*\cal{ F}$ (\resp $\R^1_{\tame}f_*\cal{ F}$)}
{$g_*\cal{ F}$ (\resp $\R^1_{\tame}g_*\cal{ F}$)}
es localmente constante y constructible, y
\ifthenelse{\boolean{orig}}
{$(\cal{ F},f)$}
{$(\cal{ F},g)$}
es cohomológicamente propio
relativamente a~$S$ en dimensión $\leq 0$ (\resp la formación de
\ifthenelse{\boolean{orig}}
{$\R^1_{\tame}f_*\cal{ F}$}
{$\R^1_{\tame}g_*\cal{ F}$}
conmuta con todo cambio de base
$S^{'}\to S$). En el caso sin paréntesis, si $\cal{ F}$ es un
haz de grupos, para toda especialización $\overline{s}_1
\to \overline{s}_2$ de puntos geométricos de~$S$, el
morfismo de especialización
$$
\ifthenelse{\boolean{orig}}
{(\R^1_{\tame}f_*\cal{ F})_{\overline{s}_2}\to
(\R^1_{\tame}f_*\cal{ F})_{\overline{s}_1}}
{(\R^1_{\tame}g_*\cal{ F})_{\overline{s}_2}\to
(\R^1_{\tame}g_*\cal{ F})_{\overline{s}_1}}
$$
es inyectivo.
\end{corollaire}
El corolario resulta inmediatamente de~
\ifthenelse{\boolean{orig}}
{\Ref{XIII.2.6}}
{\Ref{XIII.2.7}}
y
de~\Ref{XIII.1.14} (\resp del análogo de~\Ref{XIII.1.14} para
\ifthenelse{\boolean{orig}}
{el $\R^1_{\tame}f_*\cal{ F}$}
{el $\R^1_{\tame}g_*\cal{ F}$},
que se demuestra como en \loccit).

\begin{corollaire} \label{XIII.2.9}
Sea
$$
\xymatrix{ U \ar[r]^{i} \ar[d]_{g} & X \ar[ld]^{f} \\ S }
$$
un diagrama conmutativo de esquemas, en el que $U$ es el abierto
complementario en~$X$ de un divisor de cruzamientos normales
relativamente a~$S$, y $f$~un morfismo propio y liso de
presentación finita. Sea~$\LL $ el conjunto de los números primos
distintos de las características residuales de~$S$. Sea~$\cal{F}$ un
haz de~$\LL $-grupos localmente constante y constructible sobre~$U$,
moderadamente ramificado sobre~$X$ relativamente a~$S$. Entonces
\ifthenelse{\boolean{orig}}
{$\R^1 f_*\cal{ F}$}
{$\R^1 g_*\cal{ F}$}
es localmente constante y constructible, y
\ifthenelse{\boolean{orig}}
{$(\cal{ F},f)$}
{$(\cal{ F},g)$}
es cohomológicamente propio
relativamente
\marginpar{390}
a~$S$ en dimensión $\leq 1$.
\end{corollaire}
El corolario resulta de~\Ref{XIII.2.8} y del hecho de que se tiene
\ifthenelse{\boolean{orig}}
{$\R_{\tame}^1f_*\cal{ F}=\R^1f_*\cal{ F}$}
{$\R_{\tame}^1g_*\cal{ F}=\R^1g_*\cal{ F}$}
(\Ref{XIII.2.3}~b)).

\endgroup
